% ====================================================================
%  APPLICATION FORM: one per job
% ====================================================================
%  This is the ONLY file you edit for a new application. The cover
%  letter, research statement, teaching statement, writing samples and
%  combined packet all read it, so the employer's name, the position
%  and any tailored paragraph are typed once and appear everywhere.
%
%  How to use it:
%    1. make new APP=spook-street-insider-risk   (copies this file)
%    2. Fill in what you know. Leave anything else empty: {}
%       Empty fields disappear from the documents without leaving
%       stray words behind.
%    3. make all APP=spook-street-insider-risk MODE=tech
%
%  Built with this blank form, every document comes out fully generic
%  ("the advertised position at your organization"). That is useful as
%  a quick check of the generic paragraphs on their own.
% ====================================================================


% ====================================================
% 1. Employer & position
%    Fill in \company OR \uni (or both). When \company is set it is
%    used as the employer; otherwise \uni is.
% ====================================================
\newcommand{\company}{}          % e.g. Spook Street Analytics
\newcommand{\uni}{}              % e.g. University of Westacre
\newcommand{\shortname}{}        % short form for second mentions, e.g. Westacre. Empty: the full name is reused.
\newcommand{\employerarticle}{}  % "the" if the name needs it mid-sentence ("at the University of Westacre"). Usually empty for companies.
\newcommand{\companyaddress}{}   % one line under the employer in the address block, e.g. Chicago, IL
\newcommand{\position}{}         % e.g. Senior Data Scientist  -> "the Senior Data Scientist position"


% ====================================================
% 2. Optional details
%    Each one adds a short phrase only when filled:
%      \team            -> " on the Insider Risk team"
%      \department      -> " in the Department of Geography"   (type the full name)
%      \school          -> " in the School of Public Policy"
%      \project         -> " on the Cold Files project"
%      \specialization  -> " with a focus on counter-intelligence"
%      \jobref          -> " (ref. JR-2041)"
% ====================================================
\newcommand{\team}{}
\newcommand{\department}{}
\newcommand{\school}{}
\newcommand{\project}{}
\newcommand{\specialization}{}
\newcommand{\jobref}{}


% ====================================================
% 3. Addressee & greeting
%    By default each mode greets a group ("Dear Hiring Team,",
%    "Dear Members of the Search Committee,"). Use these to change it.
% ====================================================
\newcommand{\contactname}{}          % surname only, e.g. Taverner  -> "Dear Ms. Taverner,"
\newcommand{\contacttitle}{}         % e.g. Dr., Professor, Ms.   (used with \contactname)
\newcommand{\recruitmentoverride}{}  % replaces the mode's group, e.g. Talent Acquisition Team
\newcommand{\salutationoverride}{}   % replaces the whole greeting, comma included, e.g. To whom it may concern,
\newcommand{\signoffoverride}{}      % replaces "Sincerely,", e.g. Kind regards,
\newcommand{\letterdate}{}           % empty = today's date


% ====================================================
% 4. One-line hook
%    A single sentence tailored to this job. It is added to the end
%    of the opening paragraph. Empty = no sentence.
% ====================================================
\newcommand{\appsentence}{}


% ====================================================
% 5. Paragraph overrides: cover letter
%    The letter has four paragraph slots. Leave a slot empty to use
%    the generic paragraph of the chosen mode
%    (cover_letter/modes/<mode>/<slot>.tex), or write your own text
%    to replace that paragraph for this application only.
% ====================================================
\newcommand{\individualizedintro}{}       % empty -> modes/<mode>/intro.tex
\newcommand{\individualizedskills}{}      % empty -> modes/<mode>/skills.tex
\newcommand{\individualizedfit}{}         % empty -> modes/<mode>/fit.tex
\newcommand{\individualizedconclusion}{}  % empty -> modes/<mode>/conclusion.tex

% Academic modes only: one research or teaching area of the department
% you want to connect to. The academia fit paragraph switches to a
% sentence about it when this is filled.
\newcommand{\deptfocus}{}                 % e.g. Cold War intelligence history


% ====================================================
% 6. Research statement (academic applications)
% ====================================================
% An extra strand added to "Future Directions". When filled, the
% statement says "Four strands" instead of "Three" and adds it as
% the fourth item.
\newcommand{\individualizedresearchexpansion}{}


% ====================================================
% 7. Teaching statement (academic applications)
% ====================================================
\newcommand{\teachingdocname}{}           % the title the ad asks for. Empty = Teaching Statement.
                                          % e.g. Teaching Philosophy, Statement of Teaching and Mentoring
\newcommand{\teachinginterestsparagraph}{} % replaces the generic "Courses I can teach" paragraph,
                                          % e.g. when the ad asks for specific courses


% ====================================================
% 7b. Diversity and service statements (academic applications)
%     Written once and sent as is: their text does not change per
%     job. Only the line under the title follows this form, plus the
%     diversity statement's title if the ad uses a different name.
% ====================================================
\newcommand{\diversitydocname}{}          % empty = Diversity Statement. e.g. Statement of
                                          % Contributions to Diversity, Equity and Inclusion


% ====================================================
% 8. Writing samples (up to three)
%    File paths are relative to a document folder, so samples kept
%    in sample_files/ are written ../sample_files/<name>.pdf
%    Leave a file empty to skip that sample.
% ====================================================
\newcommand{\writingsampleintro}{}        % optional sentence under the contents
\newcommand{\sampleonefile}{}
\newcommand{\sampleonecitation}{}
\newcommand{\sampletwofile}{}
\newcommand{\sampletwocitation}{}
\newcommand{\samplethreefile}{}
\newcommand{\samplethreecitation}{}


% ====================================================
% 9. Documents for this application
%    yes/no for each document the job asks for. `make all` builds
%    every document marked yes and merges them, in this order and
%    with PDF bookmarks, into one Application_Packet.pdf for portals
%    that take a single upload.
% ====================================================
\newcommand{\packetletter}{yes}
\newcommand{\packetresume}{no}        % industry jobs usually ask for a resume
\newcommand{\packetcv}{no}            % academic jobs usually ask for a CV
\newcommand{\packetresearch}{no}
\newcommand{\packetteaching}{no}
\newcommand{\packetdiversity}{no}
\newcommand{\packetservice}{no}
\newcommand{\packetsamples}{no}

% Your resume and CV are made outside the kit, so give their paths here.
% Paths are relative to a document folder (../ = the kit's top folder).
% A path is only used when its switch above is yes.
\newcommand{\resumefile}{../sample_files/sample_resume.pdf}  % your resume as a PDF
\newcommand{\cvfile}{../sample_files/sample_cv.pdf}          % your CV as a PDF
